\subsection{発展的な読み物}
第13章で紹介される読み物を、日本語で読む目的に合わせて整理する。

\par\smallskip\noindent\refstepcounter{table}\label{tab:ch13-07}表 \thetable: 発展的な読み物の整理\par\smallskip
表 \thetable は、発展的な読み物の整理を比較・確認しやすい形で整理したものである。\par\smallskip
\begin{longtable}{>{\raggedright\arraybackslash}p{0.25\linewidth}>{\raggedright\arraybackslash}p{0.7\linewidth}}
\toprule
文献 & 読むと得られること \\
\midrule
\endhead
Viega, Building Secure Software & 安全なソフトウェアを作り、保守する活動の全体像を学べる。監査やサービス監視も含む。 \\
Kohnfelder, Designing Secure Software & 設計と実装の段階でのセキュリティ活動、安全なプログラミング、ウェブセキュリティを学べる。 \\
Axelrod, Engineering Safe and Secure Software Systems & リスク管理の観点から、安全でセキュアなソフトウェアシステムを工学的に扱う方法を学べる。 \\
Allen et al., Software Security Engineering & プロジェクト管理者向けに、セキュリティ工学を開発プロジェクトへ組み込む方法を学べる。 \\
Bishop, Computer Security & コンピュータセキュリティの基礎理論、アクセス制御、脆弱性、攻撃、評価を広く学べる。 \\
Dotson, Practical Cloud Security & クラウド環境の資産、権限、ネットワーク、運用上の注意点を実務的に学べる。 \\
\bottomrule
\end{longtable}
